% Shared preamble for the PrimeTensor-NS prose proof documents.
% Loaded by main.tex and by standalone.tex; never loaded by a per-module
% fragment, so that every fragment stays \input-able.
%
% ENGINE: LuaLaTeX.  This will not compile under pdfLaTeX.  The sources are
% full of Lean's Unicode -- 128 distinct non-ASCII characters across the
% tree -- and a Unicode engine lets the fragments quote it literally instead
% of transliterating it.  build.sh invokes lualatex.

\usepackage{fontspec}
\usepackage{unicode-math}
\setmainfont{Latin Modern Roman}
\setsansfont{Latin Modern Sans}
\setmonofont{DejaVu Sans Mono}[Scale=MatchLowercase]
\setmathfont{Latin Modern Math}

\usepackage[margin=1.1in]{geometry}
\usepackage{amsmath, amsthm, mathtools}   % not amssymb: unicode-math supersedes it
\usepackage{fvextra}
\usepackage{xcolor}
\usepackage{enumitem}
\usepackage[hidelinks]{hyperref}

% Lean identifiers are long, unbreakable typewriter words.  Let TeX stretch
% interword space rather than overflow the measure when one lands badly.
\tolerance=1500
\emergencystretch=3em

\providecommand{\fint}{\ensuremath{\rlap{--}\!\int}}

% ---------------------------------------------------------------- theorems
\theoremstyle{definition}
\newtheorem{definition}{Definition}[section]
\newtheorem{notation}[definition]{Notation}
\newtheorem{convention}[definition]{Convention}

\theoremstyle{plain}
\newtheorem{theorem}[definition]{Theorem}
\newtheorem{proposition}[definition]{Proposition}
\newtheorem{lemma}[definition]{Lemma}
\newtheorem{corollary}[definition]{Corollary}

\theoremstyle{remark}
\newtheorem{remark}[definition]{Remark}
\newtheorem{designnote}[definition]{Design note}

% ------------------------------------------------------- Lean cross-links
% \leanfile{PrimeTensor/Depth.lean} names the source file a section renders.
\newcommand{\leanfile}[1]{%
  \par\noindent\textit{Source:}\; \texttt{#1}\par\medskip}
\newcommand{\leanref}[1]{%
  \par\nopagebreak\vspace{-0.4em}%
  \hspace*{\fill}{\small\textsf{Lean:}\;\texttt{#1}}\par\nopagebreak}
\newcommand{\lean}[1]{\texttt{#1}}

% --------------------------------------------------------------- notation
% Generic only.  Fragments must not define macros of their own -- they are
% \input into a shared document -- and must not need any, or preamble.tex
% would churn once per module.  Write Lean-derived names with these:
%
%   \name{Axis}   a type, function or field name        ->  roman
%   \ctor{succ}   a constructor or literal              ->  sans serif
%   \lean{Foo.bar}  a Lean declaration, cited verbatim  ->  typewriter
%
\newcommand{\name}[1]{\mathrm{#1}}
\newcommand{\ctor}[1]{\mathsf{#1}}
\newcommand{\size}[1]{\lvert #1 \rvert}
\newcommand{\defeq}{\equiv}

% ------------------------------------------------------- Lean code blocks
% Lean snippets are quoted verbatim from the source, Unicode included, and
% check_fragment.py holds them to that.  The environment is fvextra's
% Verbatim and NOT listings: listings reorders characters adjacent to
% non-ASCII ones under LuaLaTeX -- it renders "(delta a b" as "delta( a b"
% and the norm bars of a quantity as both bars then the quantity -- so a
% listings block silently misquotes the very source it claims to reproduce.
% Verbatim gives up keyword colouring and gets the characters right, which
% is the trade this document needs.  fvextra (over plain fancyvrb) supplies
% breaklines, so a long Lean signature wraps instead of running off the page.
% Named leancode, not lean: \DefineVerbatimEnvironment{lean} would define
% \lean and clobber the \lean{...} macro above.
\definecolor{leanrule}{HTML}{7A9CC0}

\DefineVerbatimEnvironment{leancode}{Verbatim}{%
  fontsize=\small,
  frame=leftline,
  framerule=0.8pt,
  rulecolor=\color{leanrule},
  framesep=0.9em,
  xleftmargin=1.1em,
  breaklines=true,
  breakanywhere=true,
  breaksymbolleft={},
  breakindent=2em,
  samepage=false,
}
